\cotefile{157-2}
\scanpage{2}
Catégories de Chemins et localisation

V

% archive sheet number: 1
\scanpage{3}
\mspage{1}
Catégories \unsure{localisées} et \unsure{théorie|th.} homotopique

\uline{1}. La catégorie des chemins.

Soit $\Phi$ la catégorie suivante
\[
\left\{\begin{aligned}
\operatorname{Ob}\Phi &\simeq \mathbb Z &&\text{On note }\{\alpha\}\ \inserted{\text{ou simplement }\alpha}\ \text{l'objet qui correspond à }\alpha\in\mathbb Z\\
\operatorname{Fl}\Phi &\simeq \inserted{\underbrace{\operatorname{Ob}\Phi}_{\text{identités}}\amalg}\mathbb Z &&\text{On note }\unsure{\underline{u}|u}_{\alpha}\ \text{la flèche qui \unsure{correspond|corresp.}}\\
&&&\text{à l'entier }\alpha\in Z\\
s(u_\alpha)&=\alpha\ \text{ si }\alpha\text{ pair},\ \alpha+1 &&\text{si }\alpha\text{ impair}\\
b(u_\alpha)&=\alpha+1\ \text{ si }\alpha\text{ pair},\ \alpha-1 &&\text{si }\alpha\text{ impair}
\end{aligned}\right.
\]
Il n'y a pas à définir la composition, car si deux flèches sont composables, \unsure{elles sont} l'identité. En fait, $\Phi$ est une catégorie \uline{ordonnée}, $\{\alpha\}\leftarrow\{\beta\}$ ssi $\alpha$ pair, $\beta=\alpha+1$.

On \deleted{\unsure{pose} $\unsure{=}$} note, pour $n\in\mathbb N$
\[
\Phi_n\subset\Phi\qquad \operatorname{Ob}\Phi_n=\{\alpha\in\operatorname{Ob}\Phi\mid\deleted{|\alpha|\leq n}\inserted{-n\leq\alpha\leq+n}\}
\]
la sous-catégorie pleine de $\Phi$ dont les objets sont les $\{\alpha\}$ avec $-n\leq\alpha\leq+n$. Ainsi
\diagram{\begin{tikzcd}[ampersand replacement=\&,column sep=1.1em]
\Phi_1=\bigl\{\ \& \underset{-1}{\bullet} \& \underset{0}{\bullet}\arrow[l]\arrow[r] \& \underset{1}{\bullet} \& \bigr\}\,,
\end{tikzcd}}
la duale est
\diagram{\begin{tikzcd}[ampersand replacement=\&,column sep=1em]
\Psi_1=\Phi_1^{o}=\bigl\{ \& {-1}\arrow[r] \& 0 \& 1\arrow[l] \& \bigr\}
\end{tikzcd}}

On peut considérer $\Phi$ comme somme \unsure{amalgamée} d'une suite ordonnée \inserted{infinie (des deux côtés)} de copies de $\Phi_1$, ou aussi
\diagram{\begin{tikzcd}[ampersand replacement=\&,column sep=.48em,scale=.76,transform shape,row sep=.25em]
\cdots \& \& \smash{\overbrace{\hspace{3.8em}}^{\Phi_1}} \& \& \smash{\overbrace{\hspace{3.8em}}^{\Phi_{\unsure{1}}}} \& \& \smash{\overbrace{\hspace{3.8em}}^{\Phi_1}} \& \& \smash{\overbrace{\hspace{3.8em}}^{\Phi_1}} \& \& \smash{\overbrace{\hspace{3.8em}}^{\Phi_1}} \& \& \cdots\\
\& {} \& \underset{-4}{\cdot}\arrow[l]\arrow[r] \& \underset{-3}{\cdot} \& \underset{-2}{\cdot}\arrow[l]\arrow[r] \& \underset{-1}{\cdot} \& \underset{0}{\bullet}\arrow[l]\arrow[r] \& \underset{1}{\cdot} \& \underset{2}{\cdot}\arrow[l]\arrow[r] \& \underset{3}{\cdot} \& \underset{4}{\cdot}\arrow[l]\arrow[r] \& {}\\
\cdots \& \& \& \smash{\underbrace{\hspace{3.8em}}_{\Psi_1}} \& \& \smash{\underbrace{\hspace{3.8em}}_{\Psi_1}} \& \& \smash{\underbrace{\hspace{3.8em}}_{\Psi_1}} \& \& \smash{\underbrace{\hspace{3.8em}}_{\Psi_1}} \& \& \cdots
\end{tikzcd}}
d'une suite similaire de copies de $\Psi_1$. Les $\Phi_n$ \inserted{pour $n$ impair, sont des \unsure{$=$} sommes de copies de $\Phi_1$ mises bout à bout et les \deleted{\illegible{1}} $\Phi_n$ pour $n$ pair $>0$, \unsure{de} \inserted{\unsure{sommes} de copies de $\Psi_1$}} % slanted
forment \unsure{un|une} \deleted{\unsure{suite infinie}} \inserted{système inductif \unsure{fil.}} avec foncteurs de transition fidèles, dont $\Phi$ est la limite inductive
\[
\Phi=\varinjlim_n\Phi_n.
\]

Mais pour la suite, on va faire des $\Phi_n$ un système \uline{projectif}, en définissant pour tout $n\leq m$ une \uline{rétraction} $\underline{\unsure{f|p|\rho}}_{m,n}$ de $\Phi_m$ sur $\Phi_n$, par

% archive sheet number: 2
\scanpage{4}
\mspage{2}
\[
\left\{\begin{aligned}
\underline{\unsure{f|p|\rho}}_{m,n} &: \Phi_m\ \unsure{\text{---}|\longrightarrow}\ \Phi_n\subset\Phi_m\\
\underline{\rho}_{n,m}(\{\unsure{\alpha_{n'}|\alpha_{n}^{i}}\})&=
\begin{cases}
\alpha_i=i\in\operatorname{Ob}\Phi_n&\unsure{\text{si}\ \ldots}\quad -n\leq\unsure{i|i'}\leq n\\
\alpha_i=n&\text{si}\quad n\leq i\leq m\\
\alpha_i=-n&\text{si}\quad -m\leq i\leq-n
\end{cases}
\end{aligned}\right.
\]
On note le système projectif d'ens ordonnés ainsi obtenu par $\mathbb{\Phi}=(\Phi_n)_{n\in\mathbb N}$
\[
\mathbb{\Phi}=(\Phi_n)_{n\in\mathbb N}\quad\in\operatorname{Pro}(\mathrm{Cat})
\]
\deleted{\unsure{Pour}} Soit $X$ une catégorie, alors pour $n\in\mathbb N$
\[
\underline{\deleted{\Phi}\mathrm{Ch}}_n(X)\overset{\mathrm{déf}}{=}\underline{\mathrm{Hom}}(\Phi_n,X)
\]
est la catégorie des diagrammes de $X$ de la forme
\diagram{\begin{tikzcd}[ampersand replacement=\&,column sep=.52em,scale=.82,transform shape]
(*)\quad A_{-n}\ \cdots \& A_{-2}\arrow[r,dash,"u_{-2}","{\unsure{\text{---}|\to}}"'] \& A_{-1} \& A_0\arrow[l,"u_{-1}"']\arrow[r,"u_0"] \& A_1\arrow[r,draw=none,"{\deleted{\to}}" description] \& A_2\arrow[l,"u_1"']\arrow[r,"u_2"] \& A_3 \& \cdots\arrow[l] \& A_n
\end{tikzcd}}
avec $2n$ flèches bout à bout, \inserted{$u_i$ $(-n\leq i\leq n)$} deux flèches consécutives étant \unsure{toujours|disposées} en sens opposés, \deleted{\unsure{soit}}
\diagram{\begin{tikzcd}[ampersand replacement=\&,column sep=1em]
\text{soit} \& \cdot\arrow[r] \& \cdot \& \cdot\arrow[l] \& {,}\quad\text{soit} \& {} \& \cdot\arrow[l]\arrow[r] \& {} \& .
\end{tikzcd}}
Le foncteur canonique déduit de $\rho_{n,m}:\Phi_m\to\Phi_n$
\[
\rho^{m,n}_{\deleted{\unsure{X}}}:\underline{\deleted{\Phi}\mathrm{Ch}}_n(X)\longrightarrow\underline{\deleted{\Phi}\mathrm{Ch}}_m(Y)\qquad n\leq m
\]
un \deleted{\unsure{les}} ``foncteur de dégénérescence'', consiste à \deleted{l'opération de} \inserted{compléter} un \deleted{suite} diagramme $(*)$, par $m-n$ termes \deleted{à chaque bout,} $A_{-n}$ avec des identités à gauche, et par $m-n$ termes avec des identités à droite, donc pour $n$ pair on trouve
\diagram{\begin{tikzcd}[ampersand replacement=\&,column sep=.32em,scale=.70,transform shape,row sep=.28em]
A_{-n} \& \cdots\arrow[r,"\mathrm{id}"] \& A_{-n}\arrow[r,"\mathrm{id}"] \& A_{-n}\arrow[r] \& A_{-(n-1)} \& \cdots \& \cdots \& A_{n-1} \& A_n\arrow[l]\arrow[r,"\mathrm{id}"] \& A_n \& A_n\arrow[l,"\mathrm{id}"']\arrow[r] \& \cdots \& A_n\\
\& \smash{\underbrace{\hspace{6em}}_{m-n\ \text{termes}}} \& \& \& \& \& \& \& \& \& \smash{\underbrace{\hspace{7em}}_{m-n\ \text{termes}}}
\end{tikzcd}}
C'est un foncteur pleinement fidèle (ce qui correspond au fait que $\rho_{n,m}$ est surjectif sur les objets et sur les flèches).

% archive sheet number: 3
\scanpage{5}
\mspage{3}
On \unsure{posera|pose.}
\[
\underline{\deleted{\Phi}\mathrm{Ch}}_\infty(X)=\underline{\mathrm{Hom}}(\underset{\inserted{\text{procatégorie}}}{\mathbb{\Phi}},X)\overset{\mathrm{déf}}{=}
\deleted{\unsure{\varinjlim_n\underline{\Phi}_n(X)}}\inserted{\varinjlim\underline{\Phi}_n(X)}
\]
\inserted{ou simplement $\underline{\mathrm{Ch}}(X)$}
% slanted
\margin{L'ens des objets de cette catégorie se notera $\mathrm{Ch}_\infty(X)$ ou $\mathrm{Ch}(X)$, de \unsure{$=$} on définit les $\mathrm{Ch}_n(X)$}

On a une inclusion pleinement fidèle
\[
\underline{\deleted{\Phi}\mathrm{Ch}}_\infty(X)\lhook\joinrel\longrightarrow\underline{\mathrm{Hom}}(\Phi,X),
\]
où la deuxième catégorie est celle des \deleted{suites} diagrammes infinis (dans les deux directions)
\diagram{\begin{tikzcd}[ampersand replacement=\&,column sep=.85em]
\cdots \& A_{-2}\arrow[l]\arrow[r,"u_{-2}"] \& A_{-1} \& \deleted{\unsure{X}}A_0\arrow[l,"u_{-1}"']\arrow[r,"u_0"] \& \deleted{\unsure{X}}A_1 \& A_2\arrow[l,"u_1"']\arrow[r,"u_2"] \& \cdots
\end{tikzcd}}
et où $\underline{\mathrm{Ch}}_\infty(X)$ s'identifie à la sous-catégorie pleine formée des objets pour lesquels les flèches $u_i$ sont des identités, pour $|i|$ assez grand.

On appellera $\underline{\mathrm{Ch}}_\infty(X)$ la \uline{catégorie des chemins} \unsure{dans} $X$, ses objets sont les chemins de $X$. Ceux qui proviennent de $\Phi_n(X)$ seront appelés \uline{chemins \unsure{standart|standard} de longueur $2n$}. (Le terme ``standart'' s'applique ici à nos conventions de notations. Ainsi nous ne considérons que des chemins de longueur \unsure{pair|paire}, et pour une longueur déterminée, nous choisissons \inserted{l'une} parmi les deux \deleted{possibilités de les décrire} \inserted{espèces possibles}, p.ex. pour les chemins de longueur $2$ nous avons pris seulement
\diagram{\begin{tikzcd}[ampersand replacement=\&]
A_{-1} \& A_0\arrow[l]\arrow[r] \& A_1
\end{tikzcd}}
et non pas
\diagram{\begin{tikzcd}[ampersand replacement=\&]
A_{-1}\arrow[r] \& A_0 \& A_1\arrow[l] \& .
\end{tikzcd}}

% archive sheet number: 4
\scanpage{6}
\mspage{4}
Je me rends compte que mes notations \unsure{ici ne|ne me} sont pas commodes pour la dualité, je vais les changer, en $u_i$ pour les flèches entre objets $\{i\}$ avec $i\geq0$, \deleted{\unsure{et}} et $v_i$ la flèche symétrique de $u_i$ p.r. à l'origine. (Le groupe des autom. de $\Phi_n$ est réduit au groupe $\mathbb Z_{/2}$ à deux \unsure{élémts|éléments}, le seul automorphisme non trivial étant la symétrie $\{\deleted{\unsure{x}}i\}\longrightarrow\{-i\}$. Donc nos \deleted{\unsure{diagrammes}} chemins seront notés désormais
\diagram{\begin{tikzcd}[ampersand replacement=\&,column sep=.9em]
(*)\quad A_{-2}\arrow[r,"v_1"] \& A_{-1} \& A_0\arrow[l,"v_0"']\arrow[r,dash,"u_0"] \& A_1 \& A_2\arrow[l,"u_1"']\arrow[r] \& \cdots
\end{tikzcd}}
(Ce n'est pas non plus \unsure{enthousiasmant}\ldots)

\deleted{On a deux \unsure{mors} de la} Soit \deleted{$\unsure{\varepsilon|\mathcal E}$}
\[
\varepsilon=\{-1,+1\}
\]
l'ens ordonné \uline{discret}, réduit aux deux \unsure{élémts|éléments} $-1,+1$, et définissons des foncteurs d'inclusion (pl. fidèles)
\[
\begin{aligned}
\alpha_n:\varepsilon\longrightarrow\Phi_n\qquad &\alpha_n(-1)=-n\\
&\alpha_n(1)=n\quad,
\end{aligned}
\]
ces foncteurs sont compatibles avec les $\rho_{n,m}$, d'où
\[
\alpha_\infty:\varepsilon\longrightarrow\underset{\underset{\mathrm{Pro\,Cat}}{\rotatebox{90}{$\in$}}}{\mathbb{\Phi}}\quad,
\]
d'où pour tout $X$ un foncteur
\[
\unsure{\mathcal L|\ell|\bigl(}\qquad\underline{\mathrm{Hom}}(\mathbb{\Phi},X)\longrightarrow\underline{\mathrm{Hom}}(\varepsilon,X)\simeq X\times X
\]
i.e.

% archive sheet number: 5
\scanpage{7}
\mspage{5}
\[
\begin{aligned}
\underline{\mathrm{Ch}}(X)&\longrightarrow X\times X\\
c&\longmapsto(\sigma(x),\beta(\unsure{x|m}))\quad.
\end{aligned}
\]
Pour tout chemin $c$ dans $X$ i.e. un diagramme $(*)$ (satisfaisant à la condition de dégénérescence à l'infini), on a une source ou origine du chemin, et un but ou une \uline{extrémité}
\[
\left\{\begin{aligned}
\sigma(x)&=A_i\quad\text{pour }\deleted{|i|\ \text{grand}}\ \deleted{i\leq-N,\ N\ \text{grand}\ \unsure{(-N}}\,i\leq-\operatorname{long}(x)\\
\deleted{\unsure{\sigma}}\beta(x)&=A_i\quad\text{pour }i\geq\operatorname{long}(\unsure{x|a})
\end{aligned}\right.
\]
La fibre du \deleted{\illegible{1}} foncteur précédent en $x,y$ sera notée
\[
\underline{\mathrm{Ch}}(X;x,y)\quad\text{ou simplement }\underline{\mathrm{Ch}}(x,y)\quad,
\]
\uline{catégorie des chemins de} $x$ à $y$.

2. \deleted{\unsure{Le}} \uline{Calcul d'une cat. de fractions} $X\Sigma^{-1}$

Avant de continuer le formalisme des chemins (composition etc), je vais \unsure{tester|tâter} l'utilité de la notion de \deleted{\unsure{deux}} catégorie de chemins, par le ``calcul'' des $X\Sigma^{-1}$, où
\[
\Sigma\subset\operatorname{Fl}(X)\qquad\inserted{\Sigma\ \text{contient les identités}}
\]
est une partie \inserted{\unsure{autrement|contenant les $1$}} \uline{quelconque} de $\operatorname{Fl}(X)$. On sait que $\operatorname{Ob}\deleted{\unsure{F}}X\Sigma^{-1}=\operatorname{Ob}X$, et tout revient à définir \unsure{se|s}
\deleted{\unsure{pr}} les
\[
\operatorname{Hom}_{X\Sigma^{-1}}(x,y)\qquad x,y\in\operatorname{Ob}X,
\]

% archive sheet number: 6
\scanpage{8}
\mspage{6}
et pour $x,y,z\in\operatorname{Ob}X$, la composition
\[
\operatorname{Hom}_\Sigma(x,y)\times\unsure{H\text{--}|\operatorname{Hom}}_\Sigma(y,z)\overset{(u,v)\to vu}{\unsure{\text{---}|\longrightarrow}}{H\text{--}}_{\Sigma}(x,z),
\]
(où on met $\operatorname{Hom}_\Sigma$ au lieu de $\operatorname{Hom}_{X\Sigma^{-1}}$, pour alléger les notations).

\deleted{Je \unsure{note} d} Grâce à $\Sigma$, on définit une sous-catégorie pleine
\[
\underline{\mathrm{Ch}}_{\Sigma\deleted{\unsure{X}}}(X)\subset\underline{\mathrm{Ch}}(X),
\]
% slanted
\margin{L'ens de ses objets est noté $\mathrm{Ch}_\Sigma(X)$, ses \unsure{élémts|éléments} sont les $\Sigma$-chemins de $x$ à $y$}
par la condition sur $c\in\deleted{\unsure{\operatorname{Ob}}}\mathrm{Ch}(X)$
\[
\begin{aligned}
&\deleted{\text{Pour }\unsure{i\to i+1|i\ \text{pair}}}\\
&\text{Pour }i\text{ impair i.e.}\quad i-1\mathrel{\underset{\deleted{\unsure{\to}}}{\xleftarrow{u_i}}}i\quad (u_i\ \deleted{\unsure{\text{est dirigée}}}\ \inserted{\text{est dite}}\\
&\text{``rétrograde''}),\quad c(u_i)\in\Sigma,
\end{aligned}
\]
i.e. les \deleted{\unsure{tronçons}} \deleted{flèches} ``tronçons'' rétrogrades dans le chemin $c$ doivent être des $\in\Sigma$. On définit de même les sous-catégories pleines
\[
\underline{\mathrm{Ch}}_\Sigma(X;x,y)=\{c\in\underline{\mathrm{Ch}}(X;x,y)\mid c\in\mathrm{Ch}_\Sigma(X)\}\unsure{.}
\]
Ou encore la fibre en $(x,y)$ du foncteur source-but pour $\underline{\mathrm{Ch}}_\Sigma(X)$
\[
\underline{\mathrm{Ch}}_\Sigma(X)\longrightarrow X\times X.
\]

On définit \unsure{une} application canonique
\[
\mathrm{Ch}_\Sigma(X;x,y)\longrightarrow\operatorname{Hom}_\Sigma(x,y),
\]

% archive sheet number: 7
\scanpage{9}
\mspage{7}
\unsure{en} associant à \unsure{tt} chemin de $x$ dans $y$ le $\Sigma$-homomorphisme de $x$ \unsure{dans} $y$ obtenu en \uline{inversant} \deleted{les} (dans $X\Sigma^{-1}$) les \deleted{flèches} \inserted{tronçons} rétrogrades (ce qui est possible, \unsure{puisqu'ils|puisqu'elles} sont dans $\Sigma$), de sorte qu'on trouve dans $X\Sigma^{-1}$ un ``chemin direct''
\diagram{\begin{tikzcd}[ampersand replacement=\&,column sep=.65em,scale=.82,transform shape]
x=A_{-n}\arrow[r] \& \cdots\deleted{A}\arrow[r] \& A_{-1}\arrow[r] \& A_0\arrow[r] \& A_1\arrow[r,dash,"{\unsure{\text{---}|\to}}" description] \& A_2 \& \cdots \& A_n=A_{n+1}=\cdots=y\,,
\end{tikzcd}}
et on compose tous ces morphismes.

En fait, on trouve \unsure{aussitôt|aussi tôt} que \deleted{\unsure{pour}} l'on \deleted{\unsure{a}} \inserted{définit ainsi} un \uline{foncteur}
\[
\underline{\mathrm{Ch}}_\Sigma(X)\longrightarrow\underline{\mathrm{Fl}}(X\Sigma^{-1}),
\]
du fait que si on a un diagramme
\diagram{\begin{tikzcd}[ampersand replacement=\&]
A_{i-1}\arrow[d,"f_{i-1}"] \& A_i\arrow[l,"u_i"']\arrow[d,"f_i"] \& (i\text{ pair})\\
A'_{i-1} \& A'_i\arrow[l,"u'_i"']
\end{tikzcd}}
avec $u_i,u'_i\deleted{\unsure{-}}$ dans $\Sigma$, ce diagramme, en tant que diagramme dans $X\Sigma^{-1}$, est commutatif avec flèches horizontales des \uline{iso}, et \unsure{reste} donc commutatif en inversant les dites flèches. D'où un foncteur induit
\[
\underline{\mathrm{Ch}}_\Sigma(X;\unsure{x},y)\longrightarrow\underline{\mathrm{Fl}}(\underset{\inserted{X'}}{\underbrace{X\Sigma^{-1}}};x,y).
\]
Mais \deleted{la catégorie $\underline{\mathrm{Fl}}$} pour tte catégorie $X'$, et des objets $x,y$ de $X'$, la catégorie $\underline{\mathrm{Fl}}(X';x,y)$ est

% archive sheet number: 8
\scanpage{10}
\mspage{8}
une catégorie \uline{discrète}, ayant comme ens. d'\unsure{objets|objet} $\operatorname{Hom}_{X'}(x,y\unsure{\}|?})$. Ainsi le foncteur précédent définit en fait une application canonique
\[
\boxed{\pi_0\underline{\mathrm{Ch}}_\Sigma(X;x,y)\longrightarrow\operatorname{Hom}_\Sigma(x,y)}
\]
Je dis que cette application est toujours bijective. Pour le voir, nous allons \uline{construire} ``from scratch'' la catégorie $X\Sigma^{-1}$, en prenant comme objets ceux de $X$, et \deleted{\unsure{en posant}} en \uline{définissant}
\[
\boxed{\operatorname{Hom}_\Sigma(x,y)=\pi_0\underline{\mathrm{Ch}}_\Sigma(X;x,y)}
\]
Il faut donc définir la composition
\[
\operatorname{Hom}_\Sigma(x,y)\times\operatorname{Hom}_\Sigma(y,z)\longrightarrow\operatorname{Hom}_\Sigma(x,z).
\]

Pour ceci, on va définir des foncteurs ``composition des chemins''
\[
\begin{aligned}
&\deleted{\Phi_n\unsure{(}}\\
\underline{\mathrm{Ch}}_m(X;x,y)\times\underline{\mathrm{Ch}}_n(X;y,z)&\longrightarrow\underline{\mathrm{Ch}}_{\deleted{\unsure{m}}m+n}(X;x,z),\\
(c,\qquad c')&\longmapsto c'\circ c,
\end{aligned}
\]
mais \unsure{en|on} nous \unsure{bornant|bornons} pour simplifier au cas où $m,n$ ont la \uline{$=$ parité}, de telle façon qu'en ``mettant bout à bout'' $c$ et $c'$ \inserted{(dans cet ordre)}, on trouve bien un chemin de \unsure{type} $n+m$.

% archive sheet number: 9
\scanpage{11}
\mspage{9}
Cette \deleted{\unsure{la com}} opération est ``transposée'' du couple de foncteurs
\[
\begin{aligned}
\Phi_n&\xrightarrow{\alpha}\Phi_{n+m}& \deleted{\unsure{i\mapsto i}}\ i\mapsto(i-m)\\
\Phi_m&\overset{\beta}{\unsure{\text{---}|\longrightarrow}}\Phi_{n+m}& i\mapsto(i+n)
\end{aligned}
\]
où on aura $\alpha(\{n\})=\beta(\{-m\})=n-m$, et on aura
\[
\Phi_{n+m}\simeq\Phi_n\amalg_e\Phi_m,
\]
i.e. $\Phi_{n+m}$ est la ``somme amalgamée'' dans $(\mathrm{Cat})$, obtenue en ``recollant'' $\Phi_n$ et $\Phi_m$, par identification de l'``élément but'' de $\Phi_n$ avec ``l'élément source'' de $\Phi_m$.

Le foncteur composition \unsure{des|de} chemins induit
\[
\underline{\mathrm{Ch}}_{m,\Sigma}(x,y)\times\underline{\mathrm{Ch}}_{n,\Sigma}(y,z)\longrightarrow\underline{\mathrm{Ch}}_{m+n,\Sigma}(x,z)
\]
d'où en passant aux $\pi_0$
\[
\operatorname{Hom}_{\Sigma,\deleted{\unsure{\Sigma}}\,\unsure{m|n}}(x,y)\times\operatorname{Hom}_{\Sigma,\deleted{\unsure{\Sigma}}n}(y,z)\longrightarrow\operatorname{Hom}_{\Sigma,m+n}(x,z),
\]
où on \deleted{\unsure{pose}} dénote, pour $a,b\in\operatorname{Ob}X$, par
\[
\operatorname{Hom}_{\Sigma,n}(a,b)\overset{\mathrm{déf}}{=}\pi_0(\underline{\mathrm{Ch}}_{m,\Sigma}(\unsure{a},b))
\]
\deleted{l'image de $\pi_0(\underline{\unsure{\mathrm{Chem}|\mathrm{Ch}}}_{n,\Sigma}(a,b))$ dans $\pi_0(\underline{\mathrm{Chem}}_{\Sigma}(a,b))$. Notons ici que comme les foncteurs de transition
\[
\underline{\mathrm{Ch}}_{\unsure{\infty}\Sigma,n}(X;a,b)\longrightarrow\underline{\mathrm{Ch}}_{\Sigma,n}(X;\unsure{x|a},b),
\]
sont pleinement fidèles, l'}

% archive sheet number: 10
\scanpage{12}
\mspage{10}
Pour en déduire une application sur les $\varinjlim$
\[
\underbrace{\operatorname{Hom}_\Sigma(x,y)}_{\substack{\Vert\\\varinjlim\operatorname{Hom}_{\Sigma,n}(x,y)}}
\times\underbrace{\operatorname{Hom}_\Sigma(y,z)}_{\substack{\Vert\\\underline{\lim}\operatorname{Hom}_{\Sigma,m}(y,z)}}
\longrightarrow\underbrace{\operatorname{Hom}_{\Sigma\deleted{\unsure{+}}}(x,z)}_{\substack{\Vert\\\underline{\lim}\operatorname{Hom}_{\Sigma,\unsure{n+m}}(x,z)}},
\]
il faut prouver \deleted{la}, pour $n\leq n',m\leq m'$, la commutativité de
\diagram{\begin{tikzcd}[ampersand replacement=\&,column sep=.8em,scale=.82,transform shape]
\operatorname{Hom}_{\Sigma,\unsure{n|m}}(x,y)\times\operatorname{Hom}_{\Sigma,\unsure{m|n}}(y,z)\arrow[r]\arrow[d] \& \operatorname{Hom}_{\Sigma,m+n}(x,y)\arrow[d]\\
\operatorname{Hom}_{\Sigma,n'}(x,y)\times\operatorname{Hom}_{\Sigma,m'}(y,z)\arrow[r] \& \operatorname{Hom}_{\Sigma,m'+n'}(x,y).
\end{tikzcd}}

En d'autres termes, que pour $c:x\longrightarrow \unsure{y'|y}$ et $c':y\longrightarrow z$ deux \inserted{$\Sigma$-}chemins composables de long. $n$ et $m$ respectivement, le composé \deleted{$\unsure{c'c}$} $c'c:x\longrightarrow z$ ne change pas, ``à homotopie près'' dans $\underline{\mathrm{Ch}}_\Sigma(x,z)$, quand on remplace $c,c'$ par les chemins $\bar c,\bar c'$ de longueurs $\unsure{n',m'|m',n'}$ déduits de $c$ et $c'$ par ``dégénérescences aux bouts''.

\deleted{Ainsi, \unsure{on} \unsure{voit que le chemin}} \inserted{Plus précisément,}
\[
\overline{c'}\,\bar c\ \deleted{\unsure{\sim}}\ \overline{c'c},
\]
le chemin composé est homotope à celui déduit de $c'c$ \unsure{par} ``dégénérescences aux bouts''.

Je viens de regarder pour $n=m=1$, $n'\deleted{\unsure{+}}=m'=\unsure{2|3|?}$, et de constater qu'il n'y a pas de \unsure{flèche|flèches} \unsure{naturelle|naturelles|reliant} entre $\overline{c'}\,\bar c$ et $\overline{c'c}$ -- l'homotopie cherchée doit

% archive sheet number: 11
\scanpage{13}
\mspage{11}
se définir à l'aide d'un ``chemin'' dans $\deleted{\unsure{\operatorname{Hom}_\Sigma(}}\underline{\mathrm{Ch}}_{\Sigma,m'+n'}(X;x,z)$. Dans le cas général, on va en fait trouver une homotopie \deleted{\unsure{$k$}}
\[
h=h_{(n,m),(n',m')},
\]
``rendant commutatif'' le diagramme de foncteurs
\diagram{\begin{tikzcd}[ampersand replacement=\&,column sep=2.4em,scale=.85,transform shape,row sep=3.3em,cells={font=\scriptsize}]
A=\underline{\mathrm{Ch}}_{\Sigma,m}(X;x,y)\times\underline{\mathrm{Ch}}_{\Sigma,n}(X;y,z)
\arrow[r,"{*_{m,n}}"]\arrow[d,"{i_{m,m'}\times j_{\unsure{m|n},n'}}","{\deleted{\unsure{i_{m}}}}"']
\& \underline{\mathrm{Ch}}_{\Sigma,m+n}(X;x,z)\arrow[d,"{k_{m+n,m'+n'}}",""{name=R,inner sep=0pt}]\\
\underline{\mathrm{Ch}}_{\Sigma,m'}(X;x,y)\times\underline{\mathrm{Ch}}_{\Sigma,n'}(X;y,z)
\arrow[r,"{*_{m',n'}}",""{name=D,inner sep=0pt}]
\& \underset{\Vert}{\underline{\mathrm{Ch}}_{\Sigma,m'+n'}(X;x,z)}\arrow[from=R,to=D,bend right=17,"h"']\\
\& B
\end{tikzcd}}
où les $i_{\alpha,\beta},j_{\alpha,\beta},k_{\alpha,\beta}$ sont les foncteurs de transition
\[
\underline{\mathrm{Ch}}_{\Sigma,\alpha}(X;a,b)\,\unsure{\text{---}|\longrightarrow}\,\underline{\mathrm{Ch}}_{\Sigma,\beta}(X;a,b)
\]
$(\alpha\leq\beta)$ pour les \deleted{couples} $(x,y),(y,z)$ et $(x,z)$ (en guise de $(a,b)$) respectivement. L'homotopie $h$ sera \unsure{un|une} \deleted{suite de \unsure{flèches}} chemin dans la catégorie $\underline{\mathrm{Hom}}(A,B)$, \deleted{reliant} allant de
\[
k_{m+n,m'+n'}\circ *_{m,n}\longrightarrow *_{m',n'}\circ(i_{m,\unsure{m|m'}}\times j_{n,n'}).
\]
L'existence d'une telle homotopie signifie \deleted{\unsure{q}} ni plus ni moins que le diagramme carré ci-dessus est commutatif, en tant que diagramme dans la catégorie $(\mathrm{Cathot})$ des ``catégories à homotopie près'', dont les objets sont

% archive sheet number: 12
\scanpage{14}
\mspage{12}
les (\unsure{petites ?|petites}) catégories, et \deleted{les \unsure{fl.}} où on définit
\[
\operatorname{Hom}_{\mathrm{Cathot}}(A,B)\overset{\mathrm{déf}}{=}
\underbrace{\pi_0(\underline{\mathrm{Hom}}(A,B))}_{\substack{\text{classes d'homotopie}\\\text{de foncteurs de }A\\\text{dans }B}}
\]
C'est cette conclusion qui sera pour nous utile, et non la description explicite du chemin $h$. Cette existence admise (cf ci-dessous la preuve), on voit donc que si on regarde, \inserted{pour $a,b\in\operatorname{Ob}X$, \deleted{\unsure{$\nu\in$}}} le système des catégories $\deleted{\unsure{\mathrm{Hom}}}\underline{\mathrm{Ch}}_{\Sigma,\unsure{\nu|n}}(X;a,b)$ \inserted{$(\nu\in\mathbb N)$} comme un \uline{ind-objet de $\mathrm{Cathot}$}, qu'on \deleted{\unsure{va}} désignera par
\[
\begin{aligned}
&\deleted{\unsure{\Phi_n C_\Sigma}}\\
&\underline{\mathrm{Ch}}_\Sigma^{\mathrm{ind}}(X;a,b)\in\operatorname{Ob}\operatorname{Ind}(\mathrm{Cathot}),
\end{aligned}
\]
on aura un \deleted{application \unsure{dans la}} \inserted{\deleted{\unsure{f}} morphisme} canonique dans $\operatorname{Ind}(\mathrm{Cathot})$
\[
\underline{\mathrm{Ch}}_\Sigma^{\mathrm{ind}}(X;x,y)\times\underline{\mathrm{Ch}}_\Sigma^{\mathrm{ind}}(X;y,z)\xrightarrow{*}\underline{\mathrm{Ch}}_\Sigma^{\mathrm{ind}}(X;x,z),
\]
appelé ``\uline{\deleted{\unsure{morphisme}} \inserted{ind-morphisme} de composition des chemins}''.

On aurait aimé le déduire d'un honnête \uline{foncteur} ``composition des chemins''
\[
\deleted{\unsure{\Phi}}\underline{\mathrm{Ch}}_\Sigma(X;x,y)\mathbin{\unsure{*|\times}}\underline{\mathrm{Ch}}_\Sigma(X;y,z)\xrightarrow{?}\underline{\mathrm{Ch}}_\Sigma(X;x,z),
\]
qu'on noterait $(c,c')\longmapsto c'\circ c$, ou $\unsure{c'c|c',c}$, mais je doute qu'un tel foncteur existe,

% archive sheet number: 13
\scanpage{15}
\mspage{13}
qui par passage aux ind-objets redonne le \deleted{foncteur} morphisme précédent, ou ce qui revient au $=$, tel que pour tout $m,n\in\mathbb N$ (de $=$ parité), il existe un $N$ \inserted{$\geq m+n$} assez grand tel que l'on ait un diagramme
\diagram{\begin{tikzcd}[ampersand replacement=\&,column sep=1.9em,scale=.85,transform shape,row sep=1.1em,cells={font=\scriptsize}]
\underline{\mathrm{Ch}}_{\Sigma,m}(X;x,y)\times\underline{\mathrm{Ch}}_{\Sigma,n}(X;y,z)\arrow[r,"{*_{m,n}}"]\arrow[dd,hook]\arrow[dr] \& \underline{\mathrm{Ch}}_{\Sigma,m+n}(X;x,z)\arrow[d,hook]\\
\& \underline{\mathrm{Ch}}_{\Sigma,N}(X;x,z)\arrow[d,hook]\\
\underline{\mathrm{Ch}}_\Sigma(X;x,y)\times\underline{\mathrm{Ch}}_\Sigma(X;y,z)\arrow[r,"?"] \& \underline{\mathrm{Ch}}_\Sigma(X;x,z)
\end{tikzcd}}
et où le triangle supérieur \deleted{est} \uline{commutatif dans $\mathrm{Cathot}$}. (Il faut supposer que pour $N$ \inserted{assez} grand, le composé
\diagram{\begin{tikzcd}[ampersand replacement=\&,row sep=.5em,column sep=1.4em]
{}\arrow[d,hook] \& \\ {}\arrow[r,"?"] \& {}
\end{tikzcd}}
se factorise par $\underline{\mathrm{Ch}}_{\Sigma,N}(X;x,z)$ \ldots)

Pour voir que les foncteurs composition $*_{m,n}$ \deleted{\unsure{forment}} \inserted{définissent} un hom de ind-objets dans $(\mathrm{Cathot})$, il suffit de le prouver dans \unsure{les|le} \deleted{\unsure{deux}} cas \unsure{particuliers|particulier} $(m',n')=(m+1,n+1)$. On va faire la figure en supposant $m,n$ \uline{impairs}.

\deleted{\diagram{\begin{tikzcd}[ampersand replacement=\&,column sep=.17em,scale=.85,transform shape,row sep=.25em,cells={font=\tiny}]
\& \smash{\overbrace{\hspace{2.4em}}^{\text{dég}}} \& \& \deleted{A} \& \& \smash{\overbrace{\hspace{4em}}^{\text{dég.}}} \& \& \& \& \& \& \smash{\overbrace{\hspace{2.7em}}^{\text{dég}}}\\
A_{-n}\arrow[r,"\mathrm{id}"] \& A_{-n} \& A_{-n+1}\arrow[l]\,\cdots \& A_{n-1}\arrow[l]\arrow[r] \& \underset{\deleted{\unsure{B_{-n}}}}{A_n} \& \underset{\substack{\Vert\\B_{-n}}}{A_n}\arrow[l,"\mathrm{id}"']\arrow[r,"\mathrm{id}"] \& B_{-n} \& B_{\unsure{n+1|-n+1}}\arrow[l]\arrow[r] \& \cdots\,B_{\unsure{n-1}} \& B_{\unsure{n}}\arrow[l]\arrow[r] \& B_m \& B_m\arrow[l,"\mathrm{id}"']\\
A_{-n}\arrow[r,"\mathrm{id}"] \& A_{-n} \& A_{-n}\arrow[l,"\mathrm{id}"'] \& \deleted{\unsure{A}}A_{\unsure{n-2|-n-2}}\arrow[r,dash] \& A_{n-1}\\
\& \smash{\underbrace{\hspace{4em}}_{\text{dég}}}
\end{tikzcd}}}

Je m'aperçois \unsure{ici|est} que je me suis \unsure{coulé|couvé} stupidement, la \deleted{\unsure{conjugaison}} \inserted{composition} de deux chemins de longueur \uline{impaire} \inserted{$m,n$} \uline{\unsure{n'est}}

% archive sheet number: 14
\scanpage{16}
\mspage{14}
\uline{pas un chemin} de longueur $m+n$, i.e. un diagramme de type $\Phi_n$, mais un diagramme de type \uline{dual} $\Phi_n^o=\Psi_n$. P.ex. pour $m=n=1$
\diagram{\begin{tikzcd}[ampersand replacement=\&,column sep=.85em,row sep=.2em]
A_0 \& A_1\arrow[l]\arrow[r] \& \underset{\substack{\Vert\\B_0}}{A_2} \& B_1\arrow[l]\arrow[r] \& B_2\\
\& \smash{\underbrace{\hspace{4.8em}}_c} \& \& \smash{\underbrace{\hspace{4.8em}}_{c'}}
\end{tikzcd}}
est de type $\Phi_2^o$ ! Donc pour ne pas trop compliquer la vie, on va dans les développements précédents se borner à $m,n$ \uline{pairs}, donc des chemins de longueur au moins $4$. C'est un peu ridicule ! On est ramené à prouver la commutativité à homotopie près dans les cas où on laisse $c'$ \deleted{\unsure{et $z$ et $c'$}} fixe, et où on envoie $\underline{\mathrm{Ch}}_{\Sigma,m}\inserted{(X;x,y)}$ dans $\underline{\mathrm{Ch}}_{\Sigma,m+2}(X;x,y)$, ou inversement $c$ fixe et $\underline{\mathrm{Ch}}_{\Sigma,n}(X;y,z)$ dans $\underline{\mathrm{Ch}}_{\Sigma,n+2}$. Prenons le premier cas, et notons un chemin \inserted{$c$} sans indices négatifs (ce n'était pas une notation bien heureuse !)
\diagram{\begin{tikzcd}[ampersand replacement=\&,column sep=.8em]
A_0\arrow[r] \& A_1 \& A_2\arrow[l] \& \cdots \& {}\& A_{2m}\arrow[l]
\end{tikzcd}}
(NB \deleted{\unsure{tous nos}} à cause de la condition $\unsure{n|m}$ pair, tous nos chemins commencent désormais par une flèche de $A_0$ dans $A_1$, jamais en sens inverse) \deleted{Il \unsure{faut} comparer les deux chemins
\diagram{\begin{tikzcd}[ampersand replacement=\&,column sep=.15em,scale=.81,transform shape,cells={font=\tiny}]
A_0\arrow[r,"\mathrm{id}","{\deleted{\unsure{A}}}"'] \& A_0 \& A_0\arrow[l,"\mathrm{id}"']\arrow[r] \& A_1 \& A_2\arrow[l] \& \cdots \& A_{2m}\arrow[l]\arrow[r,"\mathrm{id}"] \& A_{2m}\arrow[r,dash,"\mathrm{id}"] \& \underset{\substack{\Vert\\\unsure{B_0}}}{A_{2m}}\arrow[r,dash,"{\unsure{\to|\leftarrow}}" description] \& B_1 \& B_2\arrow[l] \& \cdots \& B_{2n}\arrow[l]
\end{tikzcd}}}

, Je vois que l'explicitation de l'homotopie, sans offrir de difficulté, est très fastidieuse à expliquer. Je \deleted{préfère} vais essayer de donner une démonstration conceptuelle, sans ``calcul'' fastidieux

% archive sheet number: 15
\scanpage{17}
\mspage{15}

Les deux foncteurs composés
\[
f=k\circ *_{mn},\quad g=*_{m',n'}\circ(i\times j)
\]
dans le \deleted{\unsure{diag}} carré \unsure{pag|page} 11, dont il faut prouver qu'ils sont homotopes\unsure{,} dans $\underline{\mathrm{Hom}}(A,B)$, peuvent être \unsure{vu|vus} comme obtenus de la façon suivante, en termes de \deleted{\unsure{deux|deux tels}} foncteurs \deleted{\unsure{de}}
\diagram{\begin{tikzcd}[ampersand replacement=\&,column sep=large]
\text{(1)} \& \varphi,\psi:\Phi_{m'+n'} \arrow[r,shift left=.45ex] \arrow[r,shift right=.45ex] \& \Phi_{m+n}
\end{tikzcd}}
d'où deux foncteurs ``transposés''
\diagram{\begin{tikzcd}[ampersand replacement=\&,column sep=large]
\text{\unsure{(2)|(1)}} \& \underline{\mathrm{Hom}}(\Phi_{m,n},X) \arrow[r,"F",shift left=.5ex] \arrow[r,"G"',shift right=.5ex] \& \underline{\mathrm{Hom}}(\Phi_{m'+n'},X)
\end{tikzcd}}
lesquels foncteurs ``induisent'' $f$ et $g$. De façon précise, désignant par $a,\deleted{\unsure{c|b}}\,\unsure{b|f}$ les objets source et but de $\Phi_{m+n}$, par $c$ l'objet qui correspond au but du chemin $c$, \inserted{de longueur $m$,} à la source du chemin $c'$ de longueur $n$ qu'il s'agit de composer, enfin par $a_0,b_0$ les \unsure{élémts|éléments} source et et but de $\Phi_{m'+n'}$, les foncteurs envisagés (1) appliquent $a_0$ dans $a\unsure{{}^{\prime}|{}^{\scriptscriptstyle ?}}$, $b_0$ dans $b$, donc définissent \deleted{\unsure{chacun}} un diagramme commutatif
\diagram{\begin{tikzcd}[ampersand replacement=\&,column sep=large]
\& \Phi_{m'+n'} \arrow[r,"\varphi",shift left=.5ex] \arrow[r,"\psi"',shift right=.5ex] \& \Phi_{m+n}\\
\text{(3)} \& \underset{\text{cat. discrète}}{\unsure{\varepsilon|\mathcal E}=\{a_0,b_0\}}\arrow[u,hook,"\alpha'"]\arrow[ur,hook,"\alpha"] \&
\end{tikzcd}}
Ceci dit, la catégorie but $B$ pour $f,g$, savoir $\underline{\mathrm{Ch}}_{\Sigma,m'+n'}(X;x,z)$, s'interprète comme la sous-catégorie pleine de $\underline{\mathrm{Hom}}(\Phi_{m,n},X)$ \deleted{\unsure{dont les}} formée

% archive sheet number: 16
